\documentclass[11pt]{article}

\usepackage[margin=1in]{geometry}
\usepackage{amsmath}
\usepackage{amssymb}
\usepackage{booktabs}
\usepackage{enumitem}
\usepackage{hyperref}
\usepackage{xcolor}

\hypersetup{
    colorlinks=true,
    linkcolor=blue,
    urlcolor=blue
}

\setlength{\parindent}{0pt}
\setlength{\parskip}{0.45em}

\title{Equity Dispersion Volatility Swap Payoff}
\author{}
\date{}

\begin{document}

\maketitle

\section{Trade Structure}

An equity dispersion volatility swap package is usually composed of:

\begin{itemize}[leftmargin=*]
    \item long volatility swaps on a basket of single-name equities;
    \item short volatility swap on the equity index, typically SPX.
\end{itemize}

The economic exposure is therefore:

\begin{equation}
    \Pi_{\mathrm{disp}}
    =
    \sum_{j \in \mathcal{S}} \Pi_j
    -
    \Pi_I ,
\end{equation}

where $\mathcal{S}$ is the set of single-name share transactions and $I$ is the index transaction. The volatility amounts or vega notionals in the term sheet already include the intended trade sizing, so no additional weight is needed unless the model starts from a raw index notional and separately allocates it across names.

\section{Payoff Versus Valuation}

The formulas below first define the terminal payoff at maturity. A mark-to-market valuation before maturity is not the payoff itself; it is the discounted expectation of that payoff under the risk-neutral measure:

\begin{equation}
    V_t
    =
    D(t,T)
    \mathbb{E}_t^{\mathbb{Q}}
    \left[
        \Pi_{\mathrm{disp}}(T)
    \right],
\end{equation}

where $D(t,T)$ is the discount factor from maturity $T$ to valuation time $t$. In practice, the expectation is estimated by combining realized variance already observed with a market-implied estimate of the remaining variance. The projected payoff formulas in the interim valuation section should therefore be read as a practical approximation to this expectation.

\section{Final Realized Volatility}

For each underlying $j$, let $P_{j,i}$ be the official closing price or level on observation day $i$. Define the daily log return as:

\begin{equation}
    r_{j,i}
    =
    \ln \left( \frac{P_{j,i}}{P_{j,i-1}} \right).
\end{equation}

The term-sheet style final realized volatility is:

\begin{equation}
    \mathrm{FRV}_j
    =
    100
    \sqrt{
        \frac{252}{N}
        \sum_{i=1}^{N} r_{j,i}^{2}
    },
\end{equation}

where $N$ is the expected number of scheduled trading days in the full observation period. The factor $252$ annualizes the realized variance, and the multiplier $100$ expresses volatility in percentage points.

\section{Single Leg Payoff}

For a volatility swap leg $j$, define:

\begin{center}
\begin{tabular}{ll}
\toprule
Symbol & Meaning \\
\midrule
$A_j$ & Volatility amount or vega notional, in USD per volatility point \\
$K_j$ & Volatility strike, in volatility points \\
$C_j$ & Volatility cap, in volatility points, if applicable \\
$\mathrm{FRV}_j$ & Final realized volatility, in volatility points \\
\bottomrule
\end{tabular}
\end{center}

If the leg is capped, the payoff to the volatility buyer is:

\begin{equation}
    \Pi_j
    =
    A_j
    \left(
        \min(\mathrm{FRV}_j, C_j) - K_j
    \right).
\end{equation}

If the cap is not applicable, the payoff to the volatility buyer is:

\begin{equation}
    \Pi_j
    =
    A_j
    \left(
        \mathrm{FRV}_j - K_j
    \right).
\end{equation}

The payoff to the volatility seller is the negative of the buyer payoff.

\section{Dispersion Package Payoff}

For a long-single-name-volatility and short-index-volatility dispersion package, the payoff is:

\begin{equation}
    \Pi_{\mathrm{disp}}
    =
    \sum_{j \in \mathcal{S}}
    A_j
    \left(
        \min(\mathrm{FRV}_j, C_j) - K_j
    \right)
    -
    A_I
    \left(
        \min(\mathrm{FRV}_I, C_I) - K_I
    \right).
\end{equation}

This position benefits when realized single-name volatility is high relative to realized index volatility. Because index variance includes covariance terms, the package is closely related to a short correlation exposure:

\begin{equation}
    \sigma_I^2
    \approx
    \sum_j w_j^2 \sigma_j^2
    +
    2 \sum_{j<k} w_j w_k \rho_{jk} \sigma_j \sigma_k .
\end{equation}

Lower realized correlation reduces index volatility relative to the basket of single-name volatilities, which is favorable for a long dispersion position.

\section{Interim Valuation Before Maturity}

At a calculation date before final maturity, suppose $n$ log returns have already been observed. The realized contribution to annualized variance is:

\begin{equation}
    \mathrm{RealizedContribution}_j
    =
    \frac{252}{N}
    \sum_{i=1}^{n} r_{j,i}^{2}.
\end{equation}

The important point is that this term is already divided by the full-period expected observation count $N$. It should not be multiplied again by $n/N$.

Let $K_{\mathrm{var},j,\mathrm{fwd}}^2$ be the forward variance rate for the remaining observation period. A practical projected variance estimate is:

\begin{equation}
    \widehat{\sigma}_{j}^{2}
    =
    \frac{252}{N}
    \sum_{i=1}^{n} r_{j,i}^{2}
    +
    \frac{N-n}{N}
    K_{\mathrm{var},j,\mathrm{fwd}}^2 .
\end{equation}

The projected final realized volatility is then:

\begin{equation}
    \widehat{\mathrm{FRV}}_j
    =
    100 \sqrt{\widehat{\sigma}_{j}^{2}} .
\end{equation}

The interim mark-to-market approximation uses the same payoff function with $\widehat{\mathrm{FRV}}_j$ replacing the unknown final $\mathrm{FRV}_j$:

\begin{equation}
    \widehat{\Pi}_{\mathrm{disp}}
    =
    \sum_{j \in \mathcal{S}}
    A_j
    \left(
        \min(\widehat{\mathrm{FRV}}_j, C_j) - K_j
    \right)
    -
    A_I
    \left(
        \min(\widehat{\mathrm{FRV}}_I, C_I) - K_I
    \right).
\end{equation}

\section{Implied Variance Input}

The theoretically cleaner input for the remaining period is a forward variance swap rate, not simply ATM implied volatility squared. The variance swap rate is extracted from a slice of the option surface at a fixed maturity $T$ by integrating across strikes.

Let $F(T)$ be the forward price for expiry $T$, and let $Q(K,T)$ denote the undiscounted price of the out-of-the-money option with strike $K$ and maturity $T$:

\begin{equation}
    Q(K,T)
    =
    \begin{cases}
        P(K,T), & K < F(T), \\
        C(K,T), & K > F(T),
    \end{cases}
\end{equation}

where $P(K,T)$ is the put price and $C(K,T)$ is the call price. A continuous-strike approximation to the annualized variance swap rate is:

\begin{equation}
    K_{\mathrm{var}}^2(T)
    \approx
    \frac{2}{T}
    \left[
        \int_{0}^{F(T)}
        \frac{P(K,T)}{K^2}
        dK
        +
        \int_{F(T)}^{\infty}
        \frac{C(K,T)}{K^2}
        dK
    \right].
\end{equation}

Equivalently:

\begin{equation}
    K_{\mathrm{var}}^2(T)
    \approx
    \frac{2}{T}
    \int_{0}^{\infty}
    \frac{Q(K,T)}{K^2}
    dK .
\end{equation}

The integral is across strikes on one fixed-expiry $T$ slice of the volatility surface. The option surface has two dimensions, strike and maturity; this calculation freezes maturity at $T$ and integrates only in the strike direction. The division by $T$ converts the option strip's implied total variance over $[0,T]$ into an annualized variance rate.

In practice, the continuous integral is replaced by a discrete sum over listed strikes:

\begin{equation}
    K_{\mathrm{var}}^2(T)
    \approx
    \frac{2}{T}
    \sum_m
    \frac{\Delta K_m}{K_m^2}
    Q(K_m,T),
\end{equation}

where $\Delta K_m$ is the strike interval around strike $K_m$. Market implementations also handle discounting, forwards, dividends, settlement conventions, bid-ask filtering, and far-wing extrapolation. The key point is that the full smile matters; a single ATM volatility does not capture skew or tail option prices.

For a remaining period from $T_1$ to $T_2$, use total implied variance and difference it:

\begin{equation}
    K_{\mathrm{var},\mathrm{fwd}}^2(T_1,T_2)
    =
    \frac{
        T_2 K_{\mathrm{var}}^2(T_2)
        -
        T_1 K_{\mathrm{var}}^2(T_1)
    }{
        T_2 - T_1
    }.
\end{equation}

If the calculation date is today and the only remaining maturity is the final valuation date, then $T_1=0$ and the forward variance rate reduces to the variance swap rate for the final maturity slice:

\begin{equation}
    K_{\mathrm{var},\mathrm{fwd}}^2(0,T)
    =
    K_{\mathrm{var}}^2(T).
\end{equation}

If only ATM implied volatility is available, a rough proxy is:

\begin{equation}
    K_{\mathrm{var},j,\mathrm{fwd}}^2
    \approx
    \sigma_{\mathrm{ATM},j,\mathrm{fwd}}^2 .
\end{equation}

This proxy is acceptable for a quick sensitivity check, but it is not a true implied variance calculation. It will be most unreliable when the underlying has strong skew, wide wings, event risk, or sparse option strikes.

\end{document}
